% -----------------------------------------------------------------------------
% -----------------------------------------------------------------------------

最小耦合 QED 的局域相互作用只有
\[
\bar\psi\gamma^\mu\psi A_\mu .
\]
所以图上最基本的事件也只有一个：一条费米子线在某个时空点与一条光子线相连。所谓 Coulomb 交换、Compton 散射、湮灭、对产生、自能、真空极化与顶角修正，并不是分别添加的新机制；它们只是同一个局域顶角被插入不同次数、连接不同外腿以后形成的不同 Green 函数项。树图描述最低阶量子振幅，圈图则表示在顶角动量守恒之后仍然没有被外部动量唯一确定的内部四动量积分。


从 Maxwell--Dirac 作用量~\eqnrefs{eq:maxwell-dirac-action} 出发，并加入协变规范固定项；为提取相互作用顶角，把作用量分成
\begin{equation}
 S=S_0+S_I,
 \qquad
 S_I=-\qe\int\dd^4 x\,
 \bar\psi(x)\gamma^\mu\psi(x)A_\mu(x).
 \label{eq:qed-S0-SI-dp7}
\end{equation}
这里 $\qe=-e$ 是以 C 为单位的电子电荷；$S_I$ 的量纲仍为作用量。\eqnrefs{eq:qed-S0-SI-dp7} 之所以没有额外的 $c$，是因为
$S=c^{-1}\int\dd^4 x\,\Lag$ 而
$\Lag_I=-\qe c\bar\psi\gamma^\mu\psi A_\mu$，两个 $c$ 恰好相消。因此任意 QED 基本顶角都必须能追溯到这一局域相互作用项。

把一条入射电子四动量 $p$、一条流入光子四动量 $q$ 与一条出射电子四动量 $p'$ 代入这一局域顶角，时空积分给出 $(2\pi\hbar)^4\delta^{(4)}(p'-p-q)$，因此每个基本顶角满足
\begin{equation}
 p+q=p'.
 \label{eq:qed-vertex-momentum-conservation-dp68}
\end{equation}
内部线可以离壳，但顶角四动量守恒从未被放宽；这一结论直接来自顶角位置的 Fourier 积分。

\begin{figure}[htbp]
 \centering
 \begin{tikzpicture}[x=1cm,y=1cm]
 \coordinate (v) at (0,0);
 \draw[fermion] (-2.2,-0.9)--(v);
 \draw[fermion] (v)--(2.2,-0.9);
 \draw[photon] (0,2.0)--(v);
 \node[font=\small,fill=white,inner sep=.8pt] at (-1.25,-0.72) {$p,s$};
 \node[font=\small,fill=white,inner sep=.8pt] at (1.25,-0.72) {$p',s'$};
 \node[font=\small,fill=white,inner sep=.8pt,anchor=west] at (0.25,1.18) {$q,\lambda$};
 \fill[fvertex] (v);
 \node[anchor=north,font=\small,fill=white,inner sep=1.2pt] at (0,-1.16) {$p+q=p'$};
\end{tikzpicture}

 \caption{QED 基本顶角的动量标记。$p,p',q$ 均为物理四动量；图中的守恒关系来自局域相互作用的时空 Fourier 积分，而不是另外添加的图规则。}
 \label{fig:qed-elementary-vertex-dp7}
\end{figure}

相互作用图规则的生成对象不是单张图，而是归一化的时序关联函数。对任意由 Heisenberg 场组成的局域或多点算符乘积 $\mathcal O$，记 $\langle\cdots\rangle_0$ 为自由真空期望值，则相互作用绘景给出 Gell--Mann--Low 关系
\begin{equation}
 \boxed{
 \langle\Tord\mathcal O\rangle
 =\frac{
 \left\langle\Tord\left\{\mathcal O
 \exp\!\left(\ii S_I/\hbar\right)\right\}\right\rangle_0}
 {\left\langle\Tord\exp\!\left(\ii S_I/\hbar\right)\right\rangle_0}.}
 \label{eq:qed-gell-mann-low-dp7}
\end{equation}
分子中的 $\exp(\ii S_I/\hbar)$ 按顶角数展开，Wick 收缩把自由场成对连接成传播子；分母则消去与 $\mathcal O$ 完全不连接的真空泡图。因此 Feynman 图是这一关联函数微扰展开的图形化记账：传播子、顶角和组合因子均由作用量展开与 Wick 收缩确定。

对一个连通 QED 图，记 $V$ 为电子--光子基本顶角数，$I_f$ 与 $I_\gamma$ 分别为内部费米子线和内部光子线数，$L$ 为积分以后仍不能由顶角四动量守恒唯一固定的独立圈四动量数。连通图的 Euler 计数关系为
\begin{equation}
 \boxed{
 L=I_f+I_\gamma-V+1.}
 \label{eq:qed-loop-euler-dp7}
\end{equation}
树图满足 $L=0$；每增加一个独立闭合回路，通常增加一个未被外部运动学固定的四动量积分。结合每个 QED 顶角恰含两条费米子半边和一条光子半边，还可在计算任何 Dirac 迹以前判断给定外态的最低微扰阶数。若 $E_f$ 与 $E_\gamma$ 分别表示外部费米子腿和光子腿数，则半边计数与\eqnrefs{eq:qed-loop-euler-dp7} 联立给出
\begin{equation}
\boxed{
V=2L+E_f+E_\gamma-2,
\qquad
\mathcal M\sim g_{\rm em}^{V}
\sim\alpha^{L+(E_f+E_\gamma)/2-1}.}
\label{eq:qed-vertex-count-order-dp76}
\end{equation}
这里第一式给满足给定外腿和圈数所需的基本顶角数，第二式给相应振幅的最低耦合阶。作为直接代入，Compton 散射有 $E_f=2,E_\gamma=2,L=0$，所以 $V=2$；纯四光子外态在 $L=0$ 时与 QED 顶角半边计数不相容，最低连通贡献出现在 $L=1,V=4$。

给内部粒子四动量 $r^\mu$，质量为 $m$，定义
\begin{equation}
 \boxed{
 \mathcal V_r\equiv r^2-m^2c^2.}
 \label{eq:qed-virtuality-definition-dp7}
\end{equation}
$\mathcal V_r$ 的 SI 单位是动量平方。外部稳定粒子由 LSZ 放在质量壳上，故 $\mathcal V_r=0$；内部线一般 $\mathcal V_r\neq0$。传播子分母正是
$(\mathcal V_r+\ii0^+)^{-1}$。因此 ``虚粒子'' 的精确数学含义是内部 Green 函数线的四动量不必满足外部渐近态的在壳条件；它不是一个可以被探测器直接捕获、但暂时 ``违反能量守恒'' 的粒子。每个顶角的四动量仍严格守恒。


Maxwell--Dirac 作用量给出唯一基本顶角，自由二次核给出电子与光子传播子，Grassmann 奇偶性给出闭合费米环负号，顶角时空积分给出四动量守恒，LSZ 固定外部在壳态，而图的关联关系决定微扰阶数。M{\o}ller、Bhabha、Compton、对产生与一圈修正只是把这些共同规则应用于不同的外态和运动学。


